\appendix

\section{Glossary}
\label{sec:glossary}

\begin{description}[nosep, font=\normalfont\ttfamily]
  \item[align] QUA statement: the named elements (all, if none) wait for each
    other's current statement to finish. Deterministic aligns compile to
    waits; non-deterministic ones add a few cycles.
  \item[cc] clock cycle, 4\nsu\ on the OPX+. Minimum play/wait 4\,cc = 16\nsu.
  \item[claims] the roles (\code{'raman'}, \code{'imaging'}) a sequence
    declares it uses; ops on unclaimed roles are trace-time errors.
  \item[config-time parameter] compiled into the QUA config once per run;
    refused as an xvar or adjust key (\code{CONFIG\_TIME\_PARAMS}).
  \item[co-rotating frame] the feedback port's storage of each hypothesis's
    Bloch vector in the frame of its next pulse's axis, $R_z(-\phi_{j,i})\,s$;
    the frame change and the $z$-rotations merge into one phasor.
  \item[ctx] the \code{OPXShotContext} a sequence receives; every method is
    trace-time and emits QUA.
  \item[drive index] the grid point $d$ the feedback loop drives; its two
    integer IFs come from a host table, so the applied frequency is exact
    (\code{drive\_index}, \code{if\_table\_hz}).
  \item[element] a QUA logical channel on its own core; ours:
    \code{raman\_80}, \code{raman\_150}, \code{raman\_switch},
    \code{imaging\_switch}, \code{artiq\_handback}, \code{apd}.
  \item[fixed] QUA 4.28 fixed point, $[-8, 8)$, wraps modulo 16.
  \item[handoff / hand-back] ARTIQ $\to$ OPX trigger (ttl32) with RF on after
    350\nsu; OPX $\to$ ARTIQ edge (ttl41) with RF off after 7.5\us\ and blocks
    released after 15\us.
  \item[host\_data] a container filled from host-side values recorded at
    trace time, never from the OPX.
  \item[mu] ARTIQ machine unit, 1\nsu\ here; \code{seconds\_to\_mu} floors.
  \item[opx\_data\_valid] int32 mask: 1 where the OPX streamed a complete
    buffer for that shot.
  \item[ParamRef] a per-shot column handed out by \code{ctx.p}; never one
    number.
  \item[pass / block] the two level-set ops on a switch element: low = ARTIQ
    RF passes, high = blocked.
  \item[QUAM] Quantum Abstract Machine, QM's component layer over QUA; not a
    dependency.
  \item[remesh (pole return)] planned: a corrective pulse to the pole, a
    reset of every hypothesis and a finer grid (section~\ref{sec:feedback:remesh}).
  \item[reset model] the feedback port's phase model: \code{reset\_if\_phase}
    on both drives before every pulse; hypothesis phase $-\omega_j t_i$, so
    the pulse starts must be exact to the cycle (0.478 turns per cycle).
  \item[role] a lab beam channel name the generic builder speaks in; the
    channel map binds it to elements.
  \item[sequence] the OPX body of one shot, \code{@opx\_sequence}, traced
    once.
  \item[shot array] a per-shot vector shipped as one flat QUA array indexed
    \code{shot*M + i}.
  \item[shot tables] \code{ShotTables}: every scalar param per shot, in
    execution order, from one sweep of \code{compute\_derived}.
  \item[sticky] an element that holds its last analog sample / digital level
    between pulses; analog plays add.
  \item[sync element] \code{raman\_switch}: the one that executes
    \code{wait\_for\_trigger}.
  \item[trace time] when the sequence function runs on the host, once, to
    emit the program.
  \item[xvardims] \code{[len(values) per xvar]}, repeats included; the leading
    axes of every container.
\end{description}

\section{Hardware milestones M1--M4}
\label{sec:milestones}

All open on 2026-09-25. From the wiki page
\emph{OPX integration -- Quantum Machines OPX+}, the phase-1 audits, and (M4)
the remesh plan.

\paragraph{M1 -- the handshake on a scope, no atoms.}
\begin{enumerate}[nosep]
  \item Trigger (ttl32) $\to$ OPX digital 1/2 edge; digital edge $\to$ RF
    envelope at the switch AOM actually blocked (MOSFET + RF switch
    turn-off). Compare with the 350\nsu\ assumption; both beams leak if it is
    longer. If longer, raise \code{t\_opx\_handoff\_artiq\_side} to the
    measured value.
  \item ARTIQ \code{sw} event $\to$ RF envelope on at the AOM (rise of the
    chain); confirm \code{t\_opx\_handoff\_opx\_side} = 1\us\ covers it.
  \item Hand-back edge $\to$ ARTIQ RF off: read the VERBOSE slack print in
    \code{wait\_for\_quantum\_machines\_handback} for the first shots; scope
    $T_e$ $\to$ RF envelope. Then set the overlap to $t_{\rm cpu,max} + 1\us +
    1\us$ (expected $\sim$8.5--10\us\ total).
  \item Run once with no OPX job: every shot ends in \code{TriggerTimeout},
    never a false edge on ttl41.
  \item Halt a job mid-block and look at digital 1/2: do they return low?
  \item Does \code{open\_qm} accept the port-less sticky digital elements?
    If not, add the dummy analog port.
  \item Does the sticky analog hold keep the IF tone oscillating? (Expected
    yes.)
  \item Does the imaging intensity servo wind up while its switch is
    blocked?
  \item Does \code{reset\_if\_phase} on the sticky drives disturb the servo?
\end{enumerate}

\paragraph{M2 -- the OPX Rabi flop.} \code{qm\_rabi\_frequency\_opx.py} /
\code{check\_rabi\_frequency\_apd\_opx.py} against the ARTIQ-native and
manual-OPX flops at the same settings. The amplitude gap (OPX drives at the
bare dds defaults 0.337/0.324\,V versus ARTIQ's $\sqrt{0.3}\times$ default)
is closed in code since 2026-09-25: the latch plays at
\code{amp(sqrt(fraction\_power\_raman))}, the
amplitude \code{RamanBeamPair.set} writes. That equal amplitudes give equal
Rabi frequencies rests on the switch-box calibration (OPX volts = DDS
amplitude); confirm it with the two flops and re-measure
\code{t\_raman\_pi\_pulse} for the OPX drive if they differ. Measure
the OPX switch-path turn-on latency (replaces 127\nsu).

\paragraph{Timing of the feedback cycle (simulator done, hardware owed).}
With the pass-3b structure the QOP simulator gave a constant per-pulse sync
overhead of 82 cycles (328\nsu, \code{t\_opx\_feedback\_align\_overhead}; 86
cycles with discrete durations, \code{\_levels}) and 1938\,cc = 7.75\us\ of
compute for $m = 21$ against the 10\us\ budget, with zero slip on every
simulated pulse. The flat-rule structure is untimed. On hardware the finish
hook recomputes the schedule from the saved durations and compares it with
\code{t\_pulse\_start\_cc} every shot (the \code{[opx feedback] slip} line, and
\code{***} when a pulse is off by more than half a cycle). Run the open-loop
sweep first and require zero slip before trusting a closed-loop run.

\paragraph{M3 -- the APD readout in OPX units.} Calibrate \code{demod2volts}
(factor 2 convention), \code{time\_of\_flight} and the \emph{sign} (on the
simulator a $+0.16$\,V DC input integrated to raw $-0.195$) against the ARTIQ
sampler path; re-take the $S_z$ endpoints \code{v\_apd\_all\_up/down\_opx} and
the width \code{std\_photon\_fraction\_opx} through an OPX \code{measure};
\code{t\_opx\_integration\_start/len} are placeholders. Then refit the light
shift and coherence pair jointly. Only then run \code{feedback\_opx.py}.

\paragraph{M4 -- the Ramsey phase scan (prerequisite for the remesh only).}
A $\pi/2$ -- wait -- $\pi/2$ sequence on the OPX with a
\code{frame\_rotation\_2pi} applied to the Raman drive(s) before the second
pulse, scanned over the rotation. It fixes which drive's frame to rotate, the
factor 2 of the double pass, and the sign that maps a frame rotation onto the
model's azimuth -- the mapping the pole-return pulse of
section~\ref{sec:feedback:remesh} needs for its axis $\psi^*$. Pair it with the
offline unit test that the model's own rotation takes a sharp posterior's
state to the pole. Not needed for the argmax loop as it stands.

\section{File map}
\label{sec:filemap}

\begin{center}\small
\begin{tabular}{@{}lp{9.4cm}@{}}
\toprule
file & role \\
\midrule
\multicolumn{2}{@{}l}{\emph{framework (generic; no kexp / qm imports at module level)}} \\
\file{k-exp/kexp/control/opx/manager.py} & \code{OPXManager}: \code{use}, \code{on\_finish\_prepare}, \code{finish}, \code{fill\_containers}, \code{check\_live\_adjust\_conflicts} \\
\file{.../builder.py} & \code{OPXProgramBuilder.trace}: the shot loop and handshake framing; \code{\_validate} \\
\file{.../context.py} & \code{OPXShotContext}: the \code{ctx} macros, \code{\_resolve} \\
\file{.../params\_bridge.py} & \code{ShotTables}, \code{ParamRef}, \code{ParamsProxy}, \code{build\_shot\_tables}, \code{derived\_dependents} \\
\file{.../sequence.py} & \code{OPXSequence}, \code{@opx\_sequence}, registry \\
\file{.../channels.py} & \code{ChannelSpec}, \code{ChannelMap} (roles $\to$ elements, handshake windows) \\
\file{.../components.py} & \code{Sticky}, \code{Element}, \code{DigitalLine}, \code{AnalogDrive}, \code{IntegratedInput}, \code{Machine} (QUAM-shaped; generates the config, serialises for provenance) \\
\file{.../units.py} & \code{s\_to\_cc}, \code{s\_to\_ns}, \code{hz\_to\_int}, \code{demod2volts}, limits \\
\file{.../trace\_log.py} & \code{OpLog}, \code{OpRecord}: provenance for the viewer \\
\file{.../viewer.py}, \file{.../bench.py} & pulse viewer bundle; \code{OPXBench} for notebooks \\
\multicolumn{2}{@{}l}{\emph{lab wiring}} \\
\file{.../opx\_config.py} & ports, truth table, \code{kexp\_channel\_map}, \code{build\_machine}/\code{build\_opx\_config}, \code{CONFIG\_TIME\_PARAMS}, \code{make\_opx\_manager} \\
\file{k-exp/kexp/base/control.py} & \code{prep\_raman}, \code{handoff\_to\_quantum\_machines}, \code{wait\_for\_quantum\_machines\_handback} \\
\file{k-exp/kexp/base/base.py} & \code{self.opx} construction; \code{finish\_prepare}; \code{cleanup\_scan\_kernel}; \code{end} \\
\file{k-exp/kexp/config/expt\_params.py} & the \code{t\_opx\_*} block (lines 431--469) \\
\file{k-exp/kexp/config/ttl\_id.py} & ttl32 / ttl41 / ttl33 (lines 64--66) \\
\file{k-exp/kexp/config/ip.py} & \code{QM\_OPX\_IP}, \code{QM\_OPX\_CLUSTER} \\
\multicolumn{2}{@{}l}{\emph{scan / data plumbing (waxx / waxa)}} \\
\file{wax/waxx-src/waxx/base/scanner.py} & \code{scan}, \code{update\_params\_from\_xvars}, \code{write\_host\_params\_to\_kernel}, \code{step\_scan}, \code{init\_xvars} \\
\file{wax/waxa-src/waxa/base/dealer.py} & \code{plug\_in\_xvars}, \code{repeat\_xvars}, \code{shuffle\_xvars}, unshuffle \\
\file{wax/waxx-src/waxx/config/data\_vault.py} & containers, \code{set\_container\_size}, \code{put\_shot\_data} \\
\file{wax/waxx-src/waxx/base/expt.py} & \code{finish\_prepare\_wax}, \code{\_notify\_shot\_complete}, \code{end\_wax}, the END payload \\
\file{wax/waxa-src/waxa/data/data\_saver.py} & liveOD side: INIT/END writes, unshuffle, extra-file attrs \\
\file{wax/waxx-src/waxx/control/artiq/TTL.py} & \code{arm}, \code{wait\_for\_edge}, \code{clear\_input\_events} \\
\multicolumn{2}{@{}l}{\emph{sequences and experiments}} \\
\file{k-exp/kexp/experiments/opx\_sequences/rabi.py} & \code{rabi\_raman\_pulse}, \code{rabi\_raman\_apd} \\
\file{.../opx\_sequences/feedback.py} & \code{bayesian\_feedback}, \code{bayesian\_feedback\_open\_loop} (\code{make\_feedback\_sequence}, \code{feedback\_constants}/\code{FeedbackConstants}, \code{build\_feedback\_shot\_tables}, \code{\_feedback\_body}, \code{feedback\_finish}); the table helpers \code{axis\_tables}, \code{level\_matrix\_tables}, \code{level\_seed\_tables}, \code{corotating\_phase\_tables}, \code{extend\_schedule\_cc}, \code{if\_tables\_hz}, \code{applied\_frequency\_hz}, \code{sincos\_lut\_table}/\code{\_emulation}, \code{exp\_lut\_table}/\code{\_emulation}, \code{feedback\_step\_gaps\_cc}, \code{scheduled\_pulse\_starts\_s}; the module docstring holds the statement list, the numerics and the simulator tables \\
\file{k-exp/kexp/experiments/qm/check\_rabi\_frequency\_apd\_opx.py} & the reference OPX experiment (\code{scan\_kernel} structure) \\
\file{.../HF\_experiments/feedback/expt\_params\_feedback\_opx.py} & the OPX feedback params (calibration placeholders, budget, the two overheads, log-weight scale, flat-rule flag, table bits, duration levels, forced flags) \\
\file{.../HF\_experiments/feedback/feedback\_opx.py}, \file{feedback\_opx\_deterministic.py} & the port's experiments (closed loop; open-loop grid sweep) \\
\file{k-exp/kexp/base/feedback.py} & the ARTIQ posterior (\code{generate\_posterior}, lines 147--368; \code{re\_initalize\_spin\_vector}, \code{remesh\_to\_centered}, \code{maybe\_remesh}, lines 470--777) and the additive host helpers for the port \code{feedback\_grid\_omega}, \code{hypothesis\_phase\_tables}, \code{schedule\_pulse\_starts}, \code{t\_raman\_pulse\_level\_set}, \code{draw\_t\_raman\_pulse\_list\_levels} (lines 1401--1556) \\
\file{k-exp/kexp/analysis/feedback\_opx.py} & \code{FeedbackOPXReplay}, \code{FeedbackOPXReplayResult}, \code{FixedPointState}, \code{posterior\_update\_fixed\_point\_emulation}, \code{run\_shot\_fixed\_point\_emulation}, \code{true\_bloch\_step}, \code{ARGMAX\_AGREEMENT\_EXPECTED} \\
\multicolumn{2}{@{}l}{\emph{tests (126 offline items in total)}} \\
\file{k-exp/tests/test\_opx.py} (39), \file{test\_opx\_ctx.py} (24), \file{test\_opx\_config.py} (19), \file{test\_opx\_viewer.py} (8), \file{opx\_sim\_fixture.py} (8), \file{test\_feedback\_opx.py} (28) & framework, context additions, config/components/handshake map, viewer, the offline simulator fixture, the feedback port \\
\multicolumn{2}{@{}l}{\emph{reports and references}} \\
\file{k-exp/docs/opx/feedback\_port\_report.md} & the feedback port in short form: phase model, exact frequency, numerics, speed, data, placeholders \\
\file{k-exp/docs/opx/remesh\_plan.md} & the remesh-with-pole-return design (section~\ref{sec:feedback:remesh}); not implemented \\
\file{k-exp/docs/opx/opx\_assessment\_2026-09-24.md} & QUAM assessment, handoff timing as coded, waste/risk findings \\
\file{k-exp.wiki/OPX-integration---Quantum-Machines-OPX+.md} & the wiki page (timeline, milestones) \\
\file{customer-weld-ucsb/} & QM's scripts for this lab (\code{configuration.py}, \code{00d/00e/00f/00g}, \code{05\_bayesian\_feedback.py}) -- superseded, kept for comparison \\
\bottomrule
\end{tabular}
\end{center}
